% from ega4-17-fr.tex
\oldpage[IV-4]{63}

\subsection{Caractérisations des morphismes non ramifiés}
\EGAReaderTarget{ega4.c04.subsec.17-4}{subsection:IV.17.4}

\begin{theorem}[17.4.1]
\label{IV.17.4.1}
Soient $f:X\to Y$ un morphisme localement de présentation finie, $x$ un point de $X$. Les propriétés suivantes sont équivalentes:
\begin{enumerate}
  \item[a)] $f$ est non ramifié au point $x$.
  \item[b)] Le morphisme diagonal $\Delta_f:X\to X\times_Y X$ est un isomorphisme local au point $x$.
  \item[b')] Si l'on pose $\Delta_f=(\psi,\theta)$, $Z=X\times_Y X$ et $z=\psi(x)$, l'homomorphisme $\theta_z^\#:\sh{O}_{Z,z}\to\sh{O}_{X,x}$ est bijectif.
  \item[b'')] Pour tout morphisme $g:Y'\to Y$, et tout point $y'\in Y'$ au-dessus de $y=f(x)$, toute $Y'$-section $s'$ de $X'=X\times_Y Y'$ telle que $x'=s'(y')$ soit au-dessus de $x$ est un isomorphisme local au point $y'$.
  \item[c)] On a $(\Omega^1_{X/Y})_x=0$.
  \item[d)] Le $k(y)$-préschéma $f^{-1}(y)$ est non ramifié sur $k(y)$ au point $x$.
  \item[d')] Le point $x$ est isolé dans $X_y=f^{-1}(y)$ (autrement dit (Err$_{\mathrm{III}}$, 20), le morphisme $f$ est quasi-fini au point $x$) et l'anneau $\sh{O}_{X_y,x}$ est un corps, extension séparable de $k(y)$.
  \item[d'')] L'anneau $\sh{O}_{X_y,x}=\sh{O}_{X,x}/\mathfrak{m}_y\sh{O}_{X,x}$ est un corps, extension finie séparable de $k(y)$.
  \item[e)] L'anneau $\sh{O}_{X,x}$ est une $\sh{O}_{Y,y}$-algèbre formellement non ramifiée pour les topologies discrètes.
\end{enumerate}
\end{theorem}

Comme $f$ est localement de type fini, le $\sh{O}_X$-Module $\Omega^1_{X/Y}$ est de type fini (16.3.9), donc il revient au même de dire que $(\Omega^1_{X/Y})_x=0$ ou qu'il existe un voisinage ouvert $U$ de $x$ tel que $\Omega^1_{X/Y}|U=0$. Compte tenu de (17.2.1), cela prouve l'équivalence de a) et c). D'autre part, si l'on pose $A=\sh{O}_{Y,y}$, $B=\sh{O}_{X,x}$, on a $(\Omega^1_{X/Y})_x=\Omega^1_{B/A}$ (16.4.15), et l'équivalence de c) et e) résulte donc de (0, 20.7.4).

Comme d') ne fait intervenir que des propriétés du morphisme $X_y\to\Spec(k(y))$, l'équivalence de a) et d') entraînera \emph{ipso facto} celle de d) et d'). D'autre part d') et d'') sont équivalentes, car il revient au même de dire que $\sh{O}_{X_y,x}$ est une $k(y)$-algèbre \emph{finie} ou que $x$ est un point isolé de $X_y$, puisque $X_y$ est un $k(y)$-préschéma localement de type fini (I, 6.4.4).

Prouvons maintenant l'équivalence de b) et b'). On peut se limiter au cas où $Y=\Spec(R)$ et $X=\Spec(S)$ sont affines et $f$ de présentation finie; alors on a $Z=\Spec(S\otimes_R S)$ et $\Delta_f$ correspond à l'homomorphisme canonique surjectif $S\otimes_R S\to S$, dont on sait que le noyau $\mathfrak{J}$ est un idéal de type fini (0, 20.4.4). Si l'on pose $\sh{J}=\widetilde{\mathfrak{J}}$, le $\sh{O}_Z$-Module $\psi_*(\sh{O}_X)=\sh{O}_Z/\sh{J}$ est donc de présentation finie, et l'hypothèse que l'homomorphisme $\theta_z^\#: \sh{O}_{Z,z}\to(\psi_*(\sh{O}_X))_z$ est bijectif entraîne qu'en remplaçant au besoin $X$ par un voisinage ouvert de $x$, l'homomorphisme $\theta:\sh{O}_Z\to\psi_*(\sh{O}_X)$ est lui-même bijectif ($0_{\mathrm I}$, 5.2.7). Cela montre donc que b') entraîne b); la réciproque est évidente.

D'autre part, l'équivalence de b) et b'') résulte de (I, 5.3.7) sans hypothèse de finitude sur $f$: la donnée d'une $Y'$-section $s':Y'\to X'$ équivaut en effet à celle d'un

